\section{Network-Conditioned Graduation}
\label{sec:network}

The baseline becomes more useful when capability conversion and unsupported thresholds depend on surrounding productive structure. Let
\begin{equation}
\Lambda_e(\Omega_t)
=
\bar\lambda_e
+\sum_{f\neq e}a_{ef}x_{f,t}q_{f,t}
+z_e(\omega_t),
\qquad a_{ef}\ge0,
\label{eq:network-lambda}
\end{equation}
where $\omega_t$ denotes external productive accessibility and $z_e(\cdot)$ summarizes access to external demand, technology, finance, suppliers, standards, logistics, or other complements. The factor $x_{f,t}$ gates the neighbor contribution: an inactive relation does not continue to raise another relation's capability-conversion rate merely because it retains a historical stock $q_{f,t}$. Let $q_e^*(\Omega_t)$ be weakly decreasing in contemporaneously available productive complementarity in the baseline case.

This formulation captures a basic point: graduation is not purely an intrinsic property of an industry. A relation surrounded by demanding customers, capable suppliers, accessible machinery, technical standards, and finance can convert operation into durable capability faster than the nominally same activity in a thin productive field. Likewise, infrastructure or market access can lower the capability threshold required for unsupported operation.

\begin{proposition}[Complementarity weakly shortens the support horizon]
\label{prop:path-dominance}
Consider two constant environments for the same relation with the same $q_0$ and $\delta$. Suppose environment $B$ satisfies
\[
\lambda_B\ge\lambda_A,
\qquad
q_B^*\le q_A^*.
\]
If graduation is feasible in both environments, then the minimum support duration in $B$ is no greater than in $A$.
\end{proposition}

\begin{proof}
Starting from the same $q_0$, the recursions satisfy $q^B_t\ge q^A_t$ for every $t$ by induction because $\lambda_B\ge\lambda_A$. Environment $B$ also requires crossing a weakly lower threshold. Therefore any date at which $A$ has graduated is a date by which $B$ has also graduated. Hence $k_B^*\le k_A^*$.
\end{proof}

Define the time-$t$ candidate graduation frontier
\begin{equation}
\G_t
=
\left\{
 e\in\E:
\frac{\Lambda_e(\Omega_t)}{\delta_e q_e^*(\Omega_t)}>1
\right\}.
\label{eq:frontier}
\end{equation}
Within the constant local approximation, $\G_t$ is the set of relations that can in principle graduate if activated and operated. It is a proposed diagnostic, not an established development index.

\subsection{Durability, dependence, and provenance}
\label{subsec:provenance}

The capability stock, the current productive environment, and the causal provenance of learning answer different questions and should not be collapsed.

First, $q_{e,t}$ is defined as \emph{durable} capability. If operating beside a temporarily supported relation $f$ teaches relation $e$ a routine, standard, engineering practice, supplier-development method, or export competence that persists after $f$ disappears, that learning belongs in $q_e$. Its origin in a supported environment does not make it fictitious or disqualify it from graduation.

Second, if the benefit from $f$ disappears when $f$ becomes inactive---for example, because $f$ is a necessary supplier, buyer, testing facility, or logistics node---that benefit is contemporaneous dependence. It belongs in $\Omega_e$, $S_e$, $D_e$, or $q_e^*(\Omega)$, not in the durable stock $q_e$.

Third, one may separately ask how much of the observed capability stock was causally generated by policy support. For an actual support path $u$ and a no-support counterfactual $0$, define the empirical attribution object
\begin{equation}
\Delta q_e^{u}(T)
=
q_e^{u}(T)-q_e^{0}(T).
\label{eq:provenance}
\end{equation}
This is a counterfactual identification problem, not part of the definition of graduation. A relation can be network-graduated at $T$ and at the same time have $\Delta q_e^u(T)>0$: support may have caused durable learning that subsequently survives without support.

The distinction prevents two opposite errors. A snapshot test should not count a currently indispensable supported neighbor as if its contribution were durable internal capability. But neither should a causal-history test erase genuine learning merely because temporary support helped produce it.

\subsection{Coupled capability propagation}
\label{subsec:coupled}

Equation~\eqref{eq:network-lambda} makes the capability stocks jointly endogenous. To isolate the implication, fix an active configuration $x$ and an external environment $\omega$ over an interval. Let
\[
X=\operatorname{diag}(x_1,\ldots,x_n),
\qquad
D=\operatorname{diag}(\delta_1,\ldots,\delta_n),
\]
and let $A=[a_{ef}]$ have nonnegative off-diagonal entries and zero diagonal. Write $z=z(\omega)$ and $\bar\lambda=(\bar\lambda_1,\ldots,\bar\lambda_n)^\top$. Then equations~\eqref{eq:capability-law} and~\eqref{eq:network-lambda} imply
\begin{equation}
q_{t+1}
=
Mq_t+b,
\qquad
M=I-D+XAX,
\qquad
b=X(\bar\lambda+z).
\label{eq:coupled-system}
\end{equation}
Assume $0<\delta_e\le1$ for every $e$ and $\bar\lambda+z\ge0$, so $M\ge0$ and $b\ge0$.

\begin{proposition}[Finite steady state under capability amplification]
\label{prop:spectral-stability}
For the fixed-configuration system in equation~\eqref{eq:coupled-system}, if
\begin{equation}
\rho(M)<1,
\label{eq:spectral-condition}
\end{equation}
then for every finite initial state $q_0\ge0$,
\begin{equation}
q_t\longrightarrow q^\infty=(I-M)^{-1}b,
\label{eq:network-steady-state}
\end{equation}
and $q^\infty$ is the unique finite fixed point of the affine recursion. Equivalently, by the standard nonsingular $M$-matrix criterion,
\begin{equation}
\rho\!\left(D^{-1}XAX\right)<1.
\label{eq:m-matrix-condition}
\end{equation}
\end{proposition}

\begin{proof}
Iterating equation~\eqref{eq:coupled-system} yields
\[
q_t=M^tq_0+\sum_{j=0}^{t-1}M^j b.
\]
If $\rho(M)<1$, then $M^t\to0$ and the Neumann series converges:
\[
\sum_{j=0}^{\infty}M^j=(I-M)^{-1}.
\]
Hence $q_t\to(I-M)^{-1}b$. Any fixed point $\tilde q$ satisfies $(I-M)\tilde q=b$; since $\rho(M)<1$, $I-M$ is nonsingular, so the fixed point is unique. Finally,
\[
I-M=D-XAX.
\]
Because $D$ is positive diagonal and $XAX\ge0$, the standard nonsingular $M$-matrix criterion gives $D-XAX$ nonsingular $M$-matrix if and only if $\rho(D^{-1}XAX)<1$, which is equivalent to $\rho(M)<1$ for the nonnegative matrix $M$.
\end{proof}

The result does not say that a real economy literally possesses an unbounded linear capability stock whenever $\rho(M)\ge1$. It says that the unsaturated linear specification no longer has a finite globally attracting steady state. At or above the boundary, diminishing returns, congestion, finite absorptive capacity, technological ceilings, or another nonlinear mechanism must eventually enter if the model is to remain realistic.

The economic content is nonetheless useful. Capability accumulation can be amplified by reciprocal productive relations before the network is analyzed as a reproductively closed system. In the linear local approximation, the relevant object is not any single $\lambda_e/\delta_e$ but the roundabout amplification encoded in $XAX$ relative to depreciation $D$.

\subsection{Unsupported dependency closure}

The coupled-capability result concerns accumulation under a fixed configuration. Graduation requires a different test: whether the resulting productive configuration survives after extraordinary support is removed.

Consider the active configuration $x^0$ at withdrawal and hold the accumulated durable capability vector $q$ fixed during the instantaneous withdrawal test. Define a support-free survival operator
\begin{equation}
\mathcal S_0(x;q)_e
=
x_e\,
\1\{q_e\ge q_e^*(\Omega(x))\}
\1\{S_e(x;0)\ge\theta_e^S\}
\1\{D_e(x;0)\ge\theta_e^D\}.
\label{eq:survival-operator}
\end{equation}
Here $\Omega(x)$ is evaluated using the currently active configuration, so inactive relations do not continue to supply contemporaneous complements. Under the pure-complementarity assumptions, $\mathcal S_0$ is monotone and satisfies $\mathcal S_0(x;q)\le x$ componentwise. Starting from $x^0$, iterate
\begin{equation}
x^{k+1}=\mathcal S_0(x^k;q).
\end{equation}

\begin{lemma}[Finite unsupported closure]
\label{lem:closure}
The sequence $\{x^k\}$ is weakly decreasing and reaches a fixed point in at most $n$ strict-removal rounds.
\end{lemma}

\begin{proof}
By construction $x^{k+1}\le x^k$. Because $x^k\in\{0,1\}^n$, every nonstationary iteration removes at least one active relation. At most $n$ relations can be removed. Therefore the sequence stabilizes after at most $n$ strict-removal rounds.
\end{proof}

\begin{definition}[Unsupported dependency closure]
The fixed point
\begin{equation}
\C_0(x^0;q)
=
\lim_{k\to\infty}\mathcal S_0^k(x^0;q)
\label{eq:closure}
\end{equation}
is the unsupported dependency closure of the current active configuration.
\end{definition}

\begin{definition}[Network graduation]
Relation $e$ is network-graduated at withdrawal if $[\C_0(x^0;q)]_e=1$.
\end{definition}

This definition is stronger than checking $q_e\ge q_e^*$ in the supported configuration. A relation that survives only because a currently support-dependent supplier or buyer remains active is not yet network-graduated. The closure therefore separates \emph{supported coexistence} from \emph{unsupported productive configuration}.

It is equally important to state what the closure does \emph{not} test. It does not ask whether the accumulated durable stock $q_e$ would have been the same in a world where support had never existed. That is the provenance question in equation~\eqref{eq:provenance}. The closure asks whether the productive relation can stand in the present without extraordinary support after contemporaneous dependencies are recursively removed.
